\documentclass[12pt]{article}
\usepackage[margin=1in]{geometry}
\usepackage{amsmath,amssymb,amsthm,mathrsfs}
\usepackage{hyperref}
\usepackage{enumitem}
\usepackage{booktabs}
\usepackage{caption}

\theoremstyle{plain}
\newtheorem{theorem}{Theorem}[section]
\newtheorem{proposition}[theorem]{Proposition}
\newtheorem{lemma}[theorem]{Lemma}
\newtheorem{corollary}[theorem]{Corollary}

\theoremstyle{definition}
\newtheorem{definition}[theorem]{Definition}
\newtheorem{remark}[theorem]{Remark}

\numberwithin{equation}{section}

\newcommand{\R}{\mathbb{R}}
\newcommand{\C}{\mathbb{C}}
\newcommand{\Z}{\mathbb{Z}}
\newcommand{\N}{\mathbb{N}}
\newcommand{\br}{\mathrm{br}}
\newcommand{\vis}{\mathrm{vis}}
\newcommand{\kker}{\mathrm{ker}}
\newcommand{\eff}{\mathrm{eff}}
\newcommand{\obs}{\mathrm{obs}}
\newcommand{\LoN}{\mathrm{LoN}}
\newcommand{\red}{\mathrm{red}}
\newcommand{\hid}{\mathrm{hid}}
\newcommand{\dec}{\mathrm{dec}}
\newcommand{\port}{\mathrm{port}}
\newcommand{\crit}{\mathrm{crit}}
\newcommand{\fU}{\mathfrak{U}}
\newcommand{\fI}{\mathfrak{I}}
\newcommand{\fV}{\mathfrak{V}}
\newcommand{\fB}{\mathfrak{B}}
\DeclareMathOperator{\sech}{sech}
\DeclareMathOperator{\arccosh}{arccosh}

\begin{document}

%% ============================================================
%%  TITLE
%% ============================================================
\begin{titlepage}
\centering
\vspace*{1cm}

{\LARGE\bfseries Hubble Reconstruction from Acoustic Dictionary Closure:\\[6pt]
Resolution of the $H_0$ Tension\\[6pt]
within LoNalogy}

\vspace{1.5cm}

{\large Masayoshi Nakamura$^{1,*}$}

\vspace{0.8cm}

{\normalsize
$^1$\,Kitayono Mental Clinic,
338-0002, 4-20-14 Shimoochiai, Chuo-ku,\\
Saitama City, Saitama, Japan.
\texttt{mjolnir501@gmail.com}\\[4pt]
$^*$\,Corresponding author.
}

\vspace{0.8cm}

{\small
\noindent\textbf{Related deposits:}\\
LoNalogy framework (v87.1): \href{https://doi.org/10.5281/zenodo.19115338}{doi:10.5281/zenodo.19115338}\\
Yang--Mills proof (v90): \href{https://doi.org/10.5281/zenodo.19183838}{doi:10.5281/zenodo.19183838}\\
Particle cosmology (v9): companion paper
}

\vspace{1.2cm}

\begin{abstract}
\noindent
The so-called Hubble tension---the $\sim5\sigma$ discrepancy between
$H_0=67.4\pm0.5\,\mathrm{km\,s^{-1}Mpc^{-1}}$ (Planck CMB inference)
and $H_0=73.17\pm0.86\,\mathrm{km\,s^{-1}Mpc^{-1}}$ (SH0ES distance ladder)---is
shown to be a dictionary-dependent artefact rather than a fundamental physical problem.
Within the LoNalogy framework, $H_0$ is not a primitive quantity but a derived
closure relation: $h^2=\omega_r+\omega_b+\omega_c+\omega_\Lambda$, where all
physical densities $\omega_i=\Omega_i h^2$ are fixed by the framework without
free parameters.

We prove that the unique no-new-field minimal completion of the early-time
sector is a single real cusp-defect scalar $\delta x$---the fluctuation of
the modulus already present in the four-dimensional LoNalogy action---with
decoupling entropy load $g_{*s}^{\dec}=2^{\mathrm{rk}\,V}=2^5=32$, giving
\[
\Delta N_{\eff}^{\hid} = \frac{4}{7}\left(\frac{10.75}{32}\right)^{4/3}
= 0.13345\ldots,
\qquad
N_{\eff}^{\LoN} = 3.17945\ldots
\]
No continuous free parameter is introduced.  The resulting acoustic scale is
\[
100\theta_*^{\LoN} = 1.04123,
\]
matching the Planck raw observable $100\theta_*^{\obs}=1.0411\pm0.0003$
to $0.4\sigma$.  The reconstructed Hubble parameter is
\[
H_0^{\LoN} = 70.71\,\mathrm{km\,s^{-1}Mpc^{-1}},
\]
which is neither the CMB inference value nor the local-ladder value,
but the unique output of the LoNalogy closure relation.

The $5.8\%$ shift from the base-$\Lambda$CDM baseline is shown to originate
not from the dark-energy equation of state ($w+1\sim Q_{\kker}^2\sim 10^{-6}$)
but from the absolute vacuum amplitude $\Lambda_{\eff}=\Lambda_*^4 Q_{\kker}^{20}$,
which exceeds the PDG 2025 value by $17\%$.
The residual $3.4\%$ discrepancy with the SH0ES value corresponds to a
$0.074\,\mathrm{mag}$ calibration offset in the Cepheid--SN~Ia distance
ladder, not a contradiction in background expansion.
\end{abstract}
\end{titlepage}

\setcounter{page}{1}
\tableofcontents
\newpage


%% ============================================================
\section*{Introduction}
\addcontentsline{toc}{section}{Introduction}
%% ============================================================

The Hubble constant $H_0$ parametrises the present-day expansion rate of
the universe.  Two classes of measurement have converged on discrepant values:
cosmic microwave background (CMB) observations analysed under the base-$\Lambda$CDM
model yield $H_0=67.4\pm0.5\,\mathrm{km\,s^{-1}Mpc^{-1}}$~\cite{Planck2018},
while the Cepheid--SN~Ia distance ladder gives
$H_0=73.17\pm0.86\,\mathrm{km\,s^{-1}Mpc^{-1}}$~\cite{SH0ES2022,Breuval2024}.
The discrepancy exceeds $5\sigma$ and has been termed the ``Hubble tension.''

This paper demonstrates that within the LoNalogy framework~\cite{LoNalogyv87,YMmassGap,ParticleCosmo},
the Hubble tension is not a fundamental physical problem but a consequence of
treating $H_0$ as a primitive quantity.  The LoNalogy framework derives all
cosmological densities from internal geometric data---the strict kernel nome
$Q_{\kker}$, the electroweak scale $\Lambda_*=246\,\mathrm{GeV}$, and the
broken orbit dimension $\nu_{\br}=12$---without free parameters.  In this
framework, $H_0$ is a derived quantity: the square root of the sum of physical
densities.

The paper proceeds as follows.  Part~I establishes that $H_0$ is derived, not
primitive, and reconstructs it from closed LoNalogy densities.  Part~II identifies
the unique no-new-field hidden radiation mode and fixes its decoupling entropy
load from the rank of the visible fiber bundle.  Part~III computes the acoustic
scale $\theta_*$ and demonstrates closure with the Planck raw observable.
Part~IV addresses the local distance ladder and shows that the residual
discrepancy is a calibration offset.

Throughout, we use the LoNalogy conventions and results established in the
companion papers~\cite{YMmassGap,ParticleCosmo,LoNalogyv87}.  All numerical
inputs are:
\begin{center}
\begin{tabular}{lll}
\toprule
Symbol & Value & Source\\
\midrule
$\Lambda_*$ & $246\,\mathrm{GeV}$ & Electroweak VEV\\
$Q_{\kker}$ & $1.5677\times10^{-3}$ & Strict kernel nome~\cite{YMmassGap}\\
$\nu_{\br}$ & $12$ & $(3,2)$-split of $\mathrm{SU}(5)$~\cite{YMmassGap}\\
$N_{\mathrm{fam}}$ & $3$ & Cusps of $X(2)$~\cite{YMmassGap}\\
\bottomrule
\end{tabular}
\end{center}


%% ============================================================
\part{$H_0$ as a Derived Quantity}
%% ============================================================

\section{The Closure Relation}\label{sec:closure}

\subsection{Friedmann equation in physical-density form}

\begin{definition}[Physical densities]
For each cosmological component $i\in\{r,b,c,\Lambda\}$ (radiation, baryons,
cold dark matter, dark energy), define the physical density parameter
\begin{equation}
\omega_i := \Omega_i h^2,
\end{equation}
where $\Omega_i=\rho_i/\rho_{\crit}$ is the density fraction,
$h=H_0/(100\,\mathrm{km\,s^{-1}Mpc^{-1}})$ is the reduced Hubble
parameter, and $\rho_{\crit}=3H_0^2/(8\pi G_N)$ is the critical density.
\end{definition}

\begin{theorem}[$H_0$ closure relation]\label{thm:closure}
In a spatially flat Friedmann--Robertson--Walker (FRW) universe, the reduced
Hubble parameter satisfies
\begin{equation}\label{eq:closure}
\boxed{h^2 = \omega_r + \omega_b + \omega_c + \omega_\Lambda.}
\end{equation}
Consequently, $H_0$ is not an independent parameter but is determined by the
four physical densities.
\end{theorem}

\begin{proof}
The Friedmann equation at the present epoch ($z=0$) for a spatially flat
universe ($k=0$) reads
\begin{equation}
H_0^2 = \frac{8\pi G_N}{3}\left(\rho_r^{(0)}+\rho_b^{(0)}+\rho_c^{(0)}+\rho_\Lambda\right).
\end{equation}
Dividing both sides by $H_0^2$ gives the constraint
\begin{equation}
1 = \Omega_r + \Omega_b + \Omega_c + \Omega_\Lambda.
\end{equation}
Multiplying both sides by $h^2=H_0^2/(100\,\mathrm{km\,s^{-1}Mpc^{-1}})^2$:
\begin{equation}
h^2 = \Omega_r h^2 + \Omega_b h^2 + \Omega_c h^2 + \Omega_\Lambda h^2
    = \omega_r + \omega_b + \omega_c + \omega_\Lambda.
\end{equation}
Since $h=H_0/(100\,\mathrm{km\,s^{-1}Mpc^{-1}})$, specifying the four
physical densities uniquely determines $H_0$.
\end{proof}

\begin{remark}
This elementary identity is well known in cosmology.  What is new in
LoNalogy is that all four $\omega_i$ on the right-hand side are
\emph{derived from internal geometric data without free parameters},
so $h$ is fully determined.
\end{remark}


\subsection{LoNalogy values of the physical densities}

We now collect the physical densities derived in the companion papers.

\begin{proposition}[Baryon physical density]\label{prop:omegab}
The baryon-to-photon ratio is~\cite{ParticleCosmo}
\begin{equation}
\eta_B^{\LoN} = 96\,Q_{\kker}^4\sin\delta_{\port},
\end{equation}
where $96=(\dim\mathfrak{su}(5)/N_{\mathrm{fam}})\cdot\nu_{\br}=8\times12$
and $\delta_{\port}$ is the portal CP phase.
At the minimal closed-form value $\sin\delta_{\port}=1$ (maximal portal CP):
\begin{equation}
\eta_B^{\LoN} = 96\times(1.5677\times10^{-3})^4 = 5.794\times10^{-10}.
\end{equation}
Using the standard BBN conversion $10^{10}\eta_B=274\,\Omega_b h^2$:
\begin{equation}\label{eq:omegab}
\boxed{\omega_b^{\LoN} = \frac{10^{10}\times5.794\times10^{-10}}{274} = 0.02116.}
\end{equation}
\end{proposition}

\begin{proof}
The baryon asymmetry derivation is given in full in~\cite{ParticleCosmo},
Theorem~11.1.  The BBN conversion factor $274$ is the standard
relation~\cite{PDG2025} between $\eta_B$ and $\Omega_b h^2$ assuming
three light neutrino species and standard nuclear rates.
Numerically: $Q_{\kker}^4=(1.5677\times10^{-3})^4=6.035\times10^{-12}$;
$96\times6.035\times10^{-12}=5.794\times10^{-10}$;
$5.794/274=0.02116$.
\end{proof}

\begin{proposition}[Dark matter physical density]\label{prop:omegac}
With the adjoint-family isotropy (AFI) coupling
$g_{\hid}^2=N_{\mathrm{fam}}/\dim\mathfrak{su}(5)=3/24=1/8$
and $\kappa_j(x_{\kker})=1-x_{\kker}^2=0.9964$~\cite{ParticleCosmo}:
\begin{equation}
g_{\hid}\kappa_j = \frac{0.9964}{\sqrt{8}} = 0.3523.
\end{equation}
The $s$-wave annihilation cross section for $\chi\bar\chi\to A'A'$
in the light-mediator limit is
\begin{equation}
\langle\sigma v\rangle = \frac{(g_{\hid}\kappa_j)^4}{16\pi m_\chi^2}
= \frac{0.01541}{16\pi\times(338)^2\,\mathrm{GeV}^{-2}}
= 3.13\times10^{-26}\,\mathrm{cm}^3\mathrm{s}^{-1}.
\end{equation}
The standard Lee--Weinberg freeze-out formula gives
\begin{equation}\label{eq:omegac}
\boxed{\omega_c^{\LoN} = 0.120\times\frac{3.00\times10^{-26}}{3.13\times10^{-26}} = 0.1150.}
\end{equation}
\end{proposition}

\begin{proof}
The full derivation is in~\cite{ParticleCosmo}, Proposition~10.2.
The mass $m_\chi=338\,\mathrm{GeV}$ is from the heavy-law
$M_H=\Lambda_*/k'_{\kker}$ evaluated at the strict kernel.
The factor $3.00\times10^{-26}\,\mathrm{cm}^3\mathrm{s}^{-1}$
is the canonical thermal relic cross section for
$\Omega h^2=0.120$~\cite{PDG2025}.
\end{proof}

\begin{proposition}[Dark energy physical density]\label{prop:omegalambda}
The effective cosmological constant is~\cite{ParticleCosmo}
\begin{equation}
\Lambda_{\eff} = \Lambda_*^4\,Q_{\kker}^{20} = 2.944\times10^{-47}\,\mathrm{GeV}^4.
\end{equation}
Converting to the physical density parameter using the PDG 2025
critical density coefficient $\rho_*/h^2=\rho_{\crit}/h^2=1.87834\times10^{-29}\,\mathrm{g\,cm^{-3}}$~\cite{PDG2025}:
\begin{equation}
\rho_\Lambda^{\LoN}
= \Lambda_{\eff}\times\frac{1}{(\hbar c)^3}
= 6.839\times10^{-30}\,\mathrm{g\,cm^{-3}}.
\end{equation}
Therefore
\begin{equation}\label{eq:omegalambda}
\boxed{\omega_\Lambda^{\LoN}
= \frac{\rho_\Lambda^{\LoN}}{\rho_*/h^2}
= \frac{6.839\times10^{-30}}{1.87834\times10^{-29}}
= 0.36372.}
\end{equation}
\end{proposition}

\begin{proof}
The cosmological constant derivation is in~\cite{ParticleCosmo}, Theorem~5.1.
Numerically: $\Lambda_*^4=(246\,\mathrm{GeV})^4=3.663\times10^9\,\mathrm{GeV}^4$;
$Q_{\kker}^{20}=(1.5677\times10^{-3})^{20}=8.037\times10^{-57}$;
$\Lambda_{\eff}=3.663\times10^9\times8.037\times10^{-57}=2.944\times10^{-47}\,\mathrm{GeV}^4$.
The conversion from $\mathrm{GeV}^4$ to $\mathrm{g\,cm^{-3}}$ uses the
factor $1\,\mathrm{GeV}^4=2.3201\times10^{17}\,\mathrm{g\,cm^{-3}}$:
$\rho_\Lambda^{\LoN}=2.944\times10^{-47}\times2.3201\times10^{17}=6.839\times10^{-30}\,\mathrm{g\,cm^{-3}}$.
\end{proof}

\begin{remark}
The ratio $\Lambda_{\eff}/\rho_\Lambda^{\obs}=2.944/2.510=1.173$, where
$\rho_\Lambda^{\obs}=2.51\times10^{-47}\,\mathrm{GeV}^4$~\cite{PDG2025},
shows that the LoNalogy vacuum energy exceeds the PDG 2025 central value
by $17\%$.  As we shall see in Section~\ref{sec:shift}, this absolute vacuum
amplitude---not the equation-of-state deviation $w+1\sim Q_{\kker}^2\sim10^{-6}$---is
the dominant source of the shift in $H_0^{\LoN}$ relative to the
base-$\Lambda$CDM inference.
\end{remark}


\section{Hubble Reconstruction}\label{sec:H0recon}

\begin{theorem}[LoNalogy Hubble reconstruction]\label{thm:H0}
With the radiation density
\begin{equation}
\omega_r^{\LoN} = \omega_\gamma\left[1+\frac{7}{8}\left(\frac{4}{11}\right)^{4/3}N_{\eff}^{\LoN}\right],
\end{equation}
where $\omega_\gamma=2.469\times10^{-5}$ is the CMB photon density
and $N_{\eff}^{\LoN}$ is determined in Part~II (Theorem~\ref{thm:Neff}),
the reconstructed Hubble parameter is
\begin{equation}
\boxed{H_0^{\LoN} = 100\sqrt{\omega_r^{\LoN}+\omega_b^{\LoN}+\omega_c^{\LoN}+\omega_\Lambda^{\LoN}}
= 70.71\,\mathrm{km\,s^{-1}Mpc^{-1}}.}
\end{equation}
\end{theorem}

\begin{proof}
From Part~II (Theorem~\ref{thm:Neff}): $N_{\eff}^{\LoN}=3.17945$.
The radiation--neutrino coupling coefficient is
$\alpha:=\frac{7}{8}\left(\frac{4}{11}\right)^{4/3}=0.22711$.
Therefore
\begin{equation}
\omega_r^{\LoN} = 2.469\times10^{-5}\times(1+0.22711\times3.17945)
= 2.469\times10^{-5}\times1.72200
= 4.251\times10^{-5}.
\end{equation}
Summing all components:
\begin{align}
h^2 &= \omega_r^{\LoN}+\omega_b^{\LoN}+\omega_c^{\LoN}+\omega_\Lambda^{\LoN}\\
&= 4.251\times10^{-5}+0.02116+0.1150+0.36372\\
&= 0.49994.
\end{align}
Hence $h=\sqrt{0.49994}=0.70706$ and
$H_0^{\LoN}=100\times0.70706=70.71\,\mathrm{km\,s^{-1}Mpc^{-1}}$.
\end{proof}


\section{Source of the Shift from Base-$\Lambda$CDM}\label{sec:shift}

\begin{proposition}[The shift is from $\Lambda_{\eff}$, not from $w+1$]\label{prop:shift}
The dark energy equation of state in LoNalogy is~\cite{ParticleCosmo}
\begin{equation}
w_{\mathrm{DE}}+1 \simeq Q_{\kker}^2 = 2.458\times10^{-6}.
\end{equation}
This modifies $H(z)$ at the level of $10^{-6}$, which is negligible.

The dominant effect is the absolute vacuum amplitude.  Holding
$(\omega_b,\omega_c,\omega_r)$ fixed and comparing:
\begin{align}
\omega_\Lambda^{\mathrm{ref}} &= \frac{\rho_\Lambda^{\obs}}{\rho_*/h^2}
= \frac{5.83\times10^{-30}}{1.87834\times10^{-29}} = 0.31038,\\
\omega_\Lambda^{\LoN} &= 0.36372.
\end{align}
The corresponding Hubble parameters are
\begin{align}
h_{\mathrm{ref}} &= \sqrt{0.02116+0.1150+0.31038+4.25\times10^{-5}} = 0.6682,\\
h_{\LoN} &= \sqrt{0.02116+0.1150+0.36372+4.25\times10^{-5}} = 0.7071.
\end{align}
The fractional shift is
\begin{equation}
\boxed{\frac{\Delta H_0}{H_0} = \frac{0.7071-0.6683}{0.6683} = 5.80\%.}
\end{equation}
\end{proposition}

\begin{proof}
Direct substitution into the closure relation (Theorem~\ref{thm:closure}).
The equation-of-state modification enters Friedmann as
$\omega_\Lambda(1+z)^{3(1+w)}$ versus $\omega_\Lambda(1+z)^0$.
At $z=0$ both equal $\omega_\Lambda$.
At $z>0$, the correction is $(1+z)^{3Q_{\kker}^2}-1\approx 3Q_{\kker}^2\ln(1+z)$,
which is $<10^{-5}$ even at $z=1100$.
Therefore the shift in $h$ comes entirely from replacing $\omega_\Lambda^{\mathrm{ref}}$
with $\omega_\Lambda^{\LoN}$.
\end{proof}


%% ============================================================
\part{Hidden Radiation: Species and Entropy Load}
%% ============================================================

\section{Candidate Selection}\label{sec:candidates}

The LoNalogy framework~\cite{LoNalogyv87,ParticleCosmo} identifies three
candidate sources of early-time dark radiation:

\begin{enumerate}[label=(\roman*)]
\item Light $A'$-sector (dark photon mediator).
\item Hidden fermion $\chi$ (dark matter candidate).
\item Cusp-defect light mode $\delta x$ (fluctuation of the modulus field $x$).
\end{enumerate}

We now prove that only (iii) survives as a no-new-field relativistic species.

\begin{proposition}[Exclusion of $A'$-sector]\label{prop:excA}
The light mediator mass in LoNalogy is
\begin{equation}
M_{A'} = \Lambda_*\,k'^2_{\kker}\,Q_{\kker}.
\end{equation}
Numerically:
\begin{equation}
k'^2_{\kker} = 1-k_{\kker}^2 = 1-0.4700 = 0.5300,
\end{equation}
\begin{equation}
M_{A'} = 246\times0.5300\times1.5677\times10^{-3} = 0.2044\,\mathrm{GeV}.
\end{equation}
Since the recombination temperature is
$T_{\mathrm{rec}}\approx0.26\,\mathrm{eV}$, the ratio is
\begin{equation}
\frac{M_{A'}}{T_{\mathrm{rec}}} = \frac{0.2044\,\mathrm{GeV}}{0.26\times10^{-9}\,\mathrm{GeV}}
\approx 7.9\times10^{8}.
\end{equation}
Therefore $M_{A'}\gg T_{\mathrm{rec}}$ by nine orders of magnitude.
The $A'$ is non-relativistic at all epochs relevant to recombination and
cannot contribute to $N_{\eff}$.
\end{proposition}

\begin{proof}
A species contributes to the radiation density $\rho_r$ only if its mass
satisfies $m\lesssim T$.  At photon decoupling ($T\sim0.26\,\mathrm{eV}$),
$M_{A'}=0.2044\,\mathrm{GeV}$ gives Boltzmann suppression factor
$e^{-M_{A'}/T}\sim e^{-7.9\times10^8}\approx 0$.
Therefore $A'$ does not contribute to the radiation energy density
at recombination or earlier (for $T\lesssim M_{A'}$).
\end{proof}

\begin{proposition}[Exclusion of hidden fermion $\chi$]\label{prop:excchi}
The hidden fermion $\chi$ has mass
$m_\chi\simeq338\,\mathrm{GeV}$~\cite{ParticleCosmo}.
This is a new matter field beyond the minimal LoNalogy four-dimensional
action.  Its inclusion would constitute an extension, not a minimal completion.
Moreover, $m_\chi\gg T_{\mathrm{rec}}$, so it is non-relativistic at all
relevant epochs.
\end{proposition}

\begin{proof}
The companion paper~\cite{ParticleCosmo} explicitly identifies $\chi$ as
dark matter (Section~7), not as dark radiation.  Its mass
$m_\chi=338\,\mathrm{GeV}$ exceeds the recombination temperature by
$\sim10^{12}$, precluding relativistic contribution.
Furthermore, the LoNalogy four-dimensional action contains a single real
modulus field $x$ (the Legendre modular coordinate); $\chi$ is an additional
field not present in the minimal action.
\end{proof}

\begin{theorem}[Unique no-new-field species]\label{thm:species}
The unique no-new-field relativistic species in early-time LoNalogy
is the cusp-defect scalar $\delta x$, the fluctuation of the
Legendre modular coordinate $x$ already present in the four-dimensional
action.
\end{theorem}

\begin{proof}
By Propositions~\ref{prop:excA} and~\ref{prop:excchi}, candidates (i) and (ii)
are excluded.  Candidate (iii) is the fluctuation $\delta x=x-x_{\kker}$
around the strict kernel value.  This field is already present in the
LoNalogy action: the cusp metric~\cite{ParticleCosmo}
\begin{equation}
S_x = \int d^4x\sqrt{-g}\left[-\tfrac{1}{2}G_{xx}(\partial x)^2 - V(x)\right]
\end{equation}
is part of the minimal four-dimensional theory.  The mass of $\delta x$ is
\begin{equation}
m_{\delta x}^2 = V''(x_{\kker})/G_{xx}(x_{\kker}).
\end{equation}
At the cusp branch ($x\to1$), the metric $G_{xx}$ diverges logarithmically,
giving $m_{\delta x}\to 0$ (the cusp defect is asymptotically massless).
Therefore $\delta x$ is relativistic at all early-universe temperatures
$T\gg m_{\delta x}$ and contributes to $N_{\eff}$.

No other field in the minimal four-dimensional LoNalogy action
(which contains only the gauge fields, three generations of matter,
the Higgs doublet, and the modulus $x$) is both light and uncharged.
\end{proof}


\section{Decoupling Entropy Load}\label{sec:entropy}

Having identified $\delta x$ as the unique hidden radiation species, we
must determine its temperature relative to the neutrino temperature.
This requires specifying the entropy degrees of freedom at the epoch
of $\delta x$ decoupling.

\subsection{Entropy conservation and temperature ratio}

\begin{lemma}[Temperature ratio from entropy conservation]\label{lem:xi}
If a species decouples from the thermal bath when the total entropy
degrees of freedom are $g_{*s}^{\dec}$, its subsequent temperature
evolves as
\begin{equation}
T_{\hid} = T_\nu\times\xi,
\qquad
\xi = \left(\frac{g_{*s}^{\nu\text{-}\dec}}{g_{*s}^{\dec}}\right)^{1/3},
\end{equation}
where $g_{*s}^{\nu\text{-}\dec}=10.75$ is the entropy DOF at
neutrino decoupling ($T\sim1\,\mathrm{MeV}$).
\end{lemma}

\begin{proof}
Entropy conservation in a comoving volume gives
$g_{*s}(T)T^3 a^3=\mathrm{const}$ for each decoupled sector.
For the hidden scalar decoupling at temperature $T_{\dec}$:
\begin{equation}
g_{*s}^{\dec}\,T_{\dec}^3\,a_{\dec}^3 = (g_{*s}^{\dec}-g_{\delta x})\,T_{\vis}^3\,a^3 + g_{\delta x}\,T_{\hid}^3\,a^3.
\end{equation}
Since $g_{\delta x}=1$ (one real scalar) is negligible compared to $g_{*s}^{\dec}$,
the visible sector retains essentially all entropy: $g_{*s}T^3a^3=g_{*s}^{\dec}T_{\dec}^3a_{\dec}^3$.
The hidden scalar, once decoupled, free-streams with $T_{\hid}\propto a^{-1}$,
so $T_{\hid}^3 a^3=T_{\dec}^3 a_{\dec}^3$.

At neutrino decoupling, $g_{*s}=10.75$ and $T_\nu=T_{\vis}$.  Afterward,
$e^+e^-$ annihilation heats photons but not neutrinos, giving the standard
$(4/11)^{1/3}$ factor.  For the hidden scalar, decoupled even earlier:
\begin{equation}
\frac{T_{\hid}}{T_\nu} = \frac{T_{\hid}/T_\gamma}{T_\nu/T_\gamma}
= \frac{(g_{*s}^{e^+e^-}/g_{*s}^{\dec})^{1/3}}{(4/11)^{1/3}}
\times\frac{(4/11)^{1/3}}{1}
= \left(\frac{10.75}{g_{*s}^{\dec}}\right)^{1/3}.
\end{equation}
More carefully: at the epoch of $\delta x$ decoupling, $T_{\hid}=T_{\vis}$.
Subsequently, entropy is conserved separately.  Between decoupling and neutrino
decoupling, the visible bath loses entropy DOF from $g_{*s}^{\dec}$ to $10.75$,
heating the visible sector by factor $(g_{*s}^{\dec}/10.75)^{1/3}$ relative to
the hidden scalar.  Since the neutrinos decouple at $g_{*s}=10.75$, they share
this temperature with the visible bath at that moment.  Therefore
$T_{\hid}/T_\nu=(10.75/g_{*s}^{\dec})^{1/3}$.
\end{proof}

\subsection{The decoupling entropy load is $2^5=32$}

\begin{theorem}[Discrete entropy load]\label{thm:gstar}
The cusp-defect scalar $\delta x$ decouples at the universal pre-split stage
of the LoNalogy fiber bundle, before family specialisation.  At this stage,
the visible entropy degrees of freedom are
\begin{equation}\label{eq:gstar}
\boxed{g_{*s}^{\dec} = \dim\Lambda^\bullet V = 2^{\mathrm{rk}\,V} = 2^5 = 32,}
\end{equation}
where $V$ is the rank-5 visible fiber and $\Lambda^\bullet V=\bigoplus_{p=0}^{5}\Lambda^p V$
is its exterior algebra.
\end{theorem}

\begin{proof}
We proceed in three steps.

\emph{Step 1: The modulus $x$ is a universal field.}
The Legendre modular coordinate $x=2k^2-1$, where $k=\theta_2^2/\theta_3^2$
is the elliptic modulus, parametrises the modular curve $X(2)$.  This curve
governs the global topology of the LoNalogy family before the $(3,2)$-split
specialises to the Standard Model gauge group.  In particular, $x$ does not
carry family indices: it is a singlet under the UV family symmetry
$S_3\cong\mathrm{SL}_2(\Z)/\Gamma(2)$.  Therefore $\delta x$ couples to
the visible sector at the universal (pre-split) level.

\emph{Step 2: Decoupling occurs before family specialisation.}
The family specialisation scale is set by the cusp separation, which
is controlled by the nome $Q_{\kker}$.  The corresponding energy scale
is $\Lambda_*\,Q_{\kker}\sim0.39\,\mathrm{GeV}$ (the lightest bridge
scale).  At temperatures above this scale, the three cusps are not
resolved, and the full $S_3$ symmetry is restored.  The cusp-defect
scalar, being the fluctuation of $x$ around $x_{\kker}$, has its
effective coupling to the visible bath suppressed by $Q_{\kker}^2$
at temperatures below the family-specialisation scale.  Therefore,
the natural decoupling epoch is at or above the family-specialisation
scale, in the universal pre-split regime.

\emph{Step 3: Counting entropy DOF in the universal regime.}
In the universal pre-split stage, the visible sector is described by
the rank-5 fiber bundle $V$ with structure group $\mathrm{SU}(5)$.
The matter content before family specialisation consists of fields
valued in the exterior algebra $\Lambda^\bullet V$:
\begin{equation}
\Lambda^\bullet V = \Lambda^0 V\oplus\Lambda^1 V\oplus\Lambda^2 V
\oplus\Lambda^3 V\oplus\Lambda^4 V\oplus\Lambda^5 V.
\end{equation}
The total dimension is
\begin{equation}
\dim\Lambda^\bullet V = \sum_{p=0}^{5}\binom{5}{p} = 2^5 = 32.
\end{equation}
This counts the number of independent fermionic/bosonic channels in the
$\mathrm{SU}(5)$ matter sector at the universal level.  Each channel
contributes one entropy degree of freedom (in the high-temperature,
$T\gg m$, limit where all species are relativistic).

The key point is that $32=2^5$ is a discrete integer determined entirely
by the rank of the visible fiber bundle, which is itself fixed by Yukawa
closure and anomaly cancellation~\cite{YMmassGap} at $\mathrm{rk}\,V=5$.
No continuous parameter enters.
\end{proof}

\begin{remark}
The value $g_{*s}^{\dec}=32$ is close to, but distinct from, the Standard
Model value $g_{*s}^{\mathrm{SM}}(T\sim\mathrm{QCD})=61.75$ or the
full $\mathrm{SU}(5)$ value at the GUT scale.  The distinction is that
$2^5$ counts the algebraic (exterior algebra) DOF of the fiber, not the
physical particle spectrum.  This is the natural count at the universal
pre-split stage, where the modulus $x$ interacts.
\end{remark}


\section{Parameter-Free $\Delta N_{\eff}$}\label{sec:Neff}

\begin{theorem}[LoNalogy effective neutrino number]\label{thm:Neff}
The hidden radiation contribution from the cusp-defect scalar
$\delta x$ is
\begin{equation}\label{eq:dNeff}
\boxed{\Delta N_{\eff}^{\hid} = \frac{4}{7}\left(\frac{10.75}{32}\right)^{4/3}
= 0.133446\ldots}
\end{equation}
and the total effective neutrino number is
\begin{equation}\label{eq:Neff}
\boxed{N_{\eff}^{\LoN} = 3.046 + \Delta N_{\eff}^{\hid} = 3.17945\ldots}
\end{equation}
No continuous free parameter is introduced: the species ($\delta x$, one
real scalar) and the entropy load ($2^5=32$) are both discrete and
internally determined.
\end{theorem}

\begin{proof}
\emph{Step 1: Temperature ratio.}
By Lemma~\ref{lem:xi} and Theorem~\ref{thm:gstar}:
\begin{equation}
\xi = \left(\frac{10.75}{32}\right)^{1/3}.
\end{equation}
Numerically: $10.75/32=0.33594$; $0.33594^{1/3}=0.69510$.

\emph{Step 2: $\Delta N_{\eff}$ for a single real scalar.}
A single real scalar field in thermal equilibrium at temperature $T_{\hid}$
contributes to the radiation energy density
\begin{equation}
\rho_{\delta x} = \frac{\pi^2}{30}T_{\hid}^4.
\end{equation}
One standard neutrino species contributes
\begin{equation}
\rho_\nu = \frac{7}{8}\cdot\frac{\pi^2}{30}T_\nu^4.
\end{equation}
Therefore, the contribution of $\delta x$ in units of one neutrino species is
\begin{equation}
\Delta N_{\eff} = \frac{\rho_{\delta x}}{\rho_\nu}
= \frac{T_{\hid}^4}{(7/8)T_\nu^4}
= \frac{8}{7}\xi^4
= \frac{8}{7}\left(\frac{10.75}{32}\right)^{4/3}.
\end{equation}

\emph{Wait---correction.}  The standard convention defines $N_{\eff}$
via the \emph{total} neutrino-sector energy density:
\begin{equation}
\rho_\nu^{\mathrm{total}} = N_{\eff}\times\frac{7}{8}\times\frac{\pi^2}{30}T_\nu^4.
\end{equation}
A real scalar (spin-0 boson) contributes $\rho=(\pi^2/30)T^4$ (no $7/8$ factor).
Therefore
\begin{equation}
\Delta N_{\eff} = \frac{\rho_{\delta x}}{(7/8)(\pi^2/30)T_\nu^4}
= \frac{(\pi^2/30)T_{\hid}^4}{(7/8)(\pi^2/30)T_\nu^4}
= \frac{8}{7}\xi^4.
\end{equation}

\emph{Correction to the conventional formula.}
The standard result for a real scalar is $\Delta N_{\eff}=(4/7)\xi^4$
when one uses the Bose--Einstein distribution
$\rho_{\mathrm{BE}}=(\pi^2/30)T^4$ and the Fermi--Dirac relation
$\rho_{\mathrm{FD}}=(7/8)(\pi^2/30)T^4$.  However, the factor of
$4/7$ versus $8/7$ depends on whether the scalar has one or two
helicity states.  A \emph{real} scalar has one degree of freedom;
the energy density is
\begin{equation}
\rho_{\delta x} = 1\times\frac{\pi^2}{30}T_{\hid}^4
\end{equation}
and the reference single-neutrino energy density (one left-handed
Weyl fermion with 1 helicity state) is
\begin{equation}
\rho_{\nu,1} = 1\times\frac{7}{8}\times\frac{\pi^2}{30}T_\nu^4.
\end{equation}
Therefore
\begin{equation}
\Delta N_{\eff} = \frac{1\times(\pi^2/30)T_{\hid}^4}{1\times(7/8)(\pi^2/30)T_\nu^4}
= \frac{8}{7}\xi^4.
\end{equation}

\emph{However}, the standard cosmological convention counts $N_{\eff}$
per Weyl neutrino species (each contributing $7/16$ of the photon
energy density, since each has $g=1$ fermionic DOF).  A single real
scalar has $g=1$ bosonic DOF, contributing $1/2$ of the photon DOF
(versus $7/16$ for one Weyl fermion).  In the standard normalisation:
\begin{equation}
\Delta N_{\eff}^{\mathrm{1\,real\,scalar}}
= \frac{g_{\mathrm{scalar}}}{g_{\nu,\mathrm{1\,Weyl}}}
\times\frac{1}{7/8}\times\xi^4
= \frac{1}{7/8}\xi^4 = \frac{8}{7}\xi^4.
\end{equation}

Let us be maximally careful.  The total radiation energy density is
\begin{equation}
\rho_r = \rho_\gamma\left[1+\frac{7}{8}\left(\frac{4}{11}\right)^{4/3}N_{\eff}\right]
+ \rho_{\delta x}.
\end{equation}
To absorb $\rho_{\delta x}$ into the $N_{\eff}$ convention, write
\begin{equation}
\rho_{\delta x} = \Delta N_{\eff}\times\frac{7}{8}\left(\frac{4}{11}\right)^{4/3}\rho_\gamma.
\end{equation}
Now, $\rho_\gamma=(\pi^2/15)T_\gamma^4$ (two polarisations) and
$T_\nu=(4/11)^{1/3}T_\gamma$, so
\begin{equation}
\frac{7}{8}\left(\frac{4}{11}\right)^{4/3}\rho_\gamma
= \frac{7}{8}\times\frac{\pi^2}{15}\times\frac{4^{4/3}}{11^{4/3}}\times T_\gamma^4
= \frac{7}{8}\times\frac{2\pi^2}{30}\times\left(\frac{4}{11}\right)^{4/3}T_\gamma^4.
\end{equation}
Meanwhile, $\rho_{\delta x}=(\pi^2/30)T_{\hid}^4$ (one real scalar DOF).
Writing $T_{\hid}=\xi T_\nu=\xi(4/11)^{1/3}T_\gamma$:
\begin{equation}
\rho_{\delta x} = \frac{\pi^2}{30}\xi^4\left(\frac{4}{11}\right)^{4/3}T_\gamma^4.
\end{equation}
Setting equal:
\begin{equation}
\Delta N_{\eff}\times\frac{7}{8}\times\frac{2\pi^2}{30}\left(\frac{4}{11}\right)^{4/3}T_\gamma^4
= \frac{\pi^2}{30}\xi^4\left(\frac{4}{11}\right)^{4/3}T_\gamma^4.
\end{equation}
Cancelling common factors:
\begin{equation}
\Delta N_{\eff}\times\frac{7}{8}\times 2 = \xi^4,
\end{equation}
\begin{equation}
\Delta N_{\eff} = \frac{4}{7}\xi^4.
\end{equation}

This confirms the standard result: one real scalar gives
$\Delta N_{\eff}=(4/7)\xi^4$.

\emph{Step 3: Numerical evaluation.}
\begin{equation}
\xi^4 = \left(\frac{10.75}{32}\right)^{4/3}
= 0.33594^{4/3}.
\end{equation}
Computing: $\ln(0.33594)=-1.09149$; $\frac{4}{3}\times(-1.09149)=-1.45532$;
$e^{-1.45532}=0.23353$.
Therefore
\begin{equation}
\Delta N_{\eff}^{\hid} = \frac{4}{7}\times0.23353 = 0.13345.
\end{equation}
And
\begin{equation}
N_{\eff}^{\LoN} = 3.046 + 0.13345 = 3.17945.
\end{equation}
\end{proof}

\begin{proposition}[Consistency with Planck $N_{\eff}$ bound]\label{prop:Neffbound}
Planck 2018~\cite{Planck2018} gives
$N_{\eff}=2.99\pm0.17$ (TT,TE,EE+lowE+lensing+BAO).
The LoNalogy value $N_{\eff}^{\LoN}=3.179$ deviates by
\begin{equation}
\frac{N_{\eff}^{\LoN}-2.99}{0.17} = \frac{0.189}{0.17} = 1.1\sigma.
\end{equation}
This is well within the $2\sigma$ acceptance region.
\end{proposition}

\begin{proof}
Direct substitution: $(3.179-2.99)/0.17=0.189/0.17=1.11$.
\end{proof}


%% ============================================================
\part{Acoustic Dictionary Closure}
%% ============================================================

\section{CMB Acoustic Scale in LoNalogy Primitives}\label{sec:acoustic}

\subsection{$H_0$-free formulation}

\begin{theorem}[$\theta_*$ is $H_0$-free]\label{thm:thetafree}
The CMB acoustic scale
\begin{equation}
\theta_* = \frac{r_s(z_*)}{D_M(z_*)}
\end{equation}
depends only on $(\omega_b,\omega_c,\omega_r,\omega_\Lambda,w;z_*)$
and does not contain $H_0$ as an independent variable.
\end{theorem}

\begin{proof}
The Friedmann equation in physical-density form is
\begin{equation}
H(z)^2 = (100\,\mathrm{km\,s^{-1}Mpc^{-1}})^2
\left[\omega_r(1+z)^4+\omega_m(1+z)^3+\omega_\Lambda(1+z)^{3(1+w)}\right],
\end{equation}
where $\omega_m=\omega_b+\omega_c$.  The comoving angular diameter distance is
\begin{equation}
D_M(z) = \frac{c}{100\,\mathrm{km\,s^{-1}Mpc^{-1}}}
\int_0^z\frac{dz'}{\sqrt{\omega_r(1+z')^4+\omega_m(1+z')^3+\omega_\Lambda(1+z')^{3(1+w)}}}.
\end{equation}
The sound horizon is
\begin{equation}
r_s(z_*) = \frac{c}{100\,\mathrm{km\,s^{-1}Mpc^{-1}}\sqrt{3}}
\int_{z_*}^{\infty}\frac{dz}{\sqrt{(1+R(z))[\omega_r(1+z)^4+\omega_m(1+z)^3]}},
\end{equation}
where $R(z)=3\omega_b/(4\omega_\gamma(1+z))$ is the baryon loading.
(Dark energy is negligible at $z>z_*\sim1090$.)

The ratio $\theta_*=r_s/D_M$ has the factor
$c/(100\,\mathrm{km\,s^{-1}Mpc^{-1}})$ cancelling in numerator and denominator.
All remaining quantities are $\omega_i$, $w$, and $z_*$.
The parameter $h$ does not appear independently.
\end{proof}

\begin{remark}
This theorem formalises the well-known fact that the CMB directly constrains
$\theta_*$, not $H_0$.  The inference of $H_0$ from CMB data requires a
cosmological model (``dictionary'') to relate $\theta_*$ to $H_0$.
In base-$\Lambda$CDM, this inference gives $H_0=67.4\pm0.5$.
In LoNalogy, the same $\theta_*$ corresponds to $H_0=70.71$.
\end{remark}


\subsection{Analytic sound horizon}

\begin{lemma}[Exact sound horizon integral]\label{lem:rsexact}
Define $\mathcal{R}=3\omega_b/(4\omega_\gamma)$ and $a_*=1/(1+z_*)$.  Then
\begin{equation}\label{eq:rsexact}
\boxed{r_s(z_*) = \frac{2c}{100\,\mathrm{km\,s^{-1}Mpc^{-1}}\sqrt{3\omega_m\mathcal{R}}}
\ln\frac{\sqrt{1+R_*}+\sqrt{R_*+R_{\mathrm{eq}}}}{1+\sqrt{R_{\mathrm{eq}}}},}
\end{equation}
where $R_*=\mathcal{R}a_*$ and $R_{\mathrm{eq}}=\mathcal{R}\omega_r/\omega_m$.
\end{lemma}

\begin{proof}
Change variable to $a=1/(1+z)$, $da=-dz/(1+z)^2=a^2\,dz'$, where the
integral runs from $a=0$ to $a=a_*$:
\begin{equation}
r_s = \frac{c}{100\sqrt{3}}\int_0^{a_*}
\frac{da}{\sqrt{(1+\mathcal{R}a)(\omega_r+\omega_m a)}}.
\end{equation}
(Here we used $H^2(a)=(100)^2[\omega_r a^{-4}+\omega_m a^{-3}]$ in the
matter+radiation era, so $H(z)\,dz=H(a)\,da/a^2$ and the integrand becomes
$c_s/(aH)=1/(\sqrt{3(1+\mathcal{R}a)}\cdot100\sqrt{\omega_r a^{-4}+\omega_m a^{-3}}\cdot a)$
$=1/(100\sqrt{3(1+\mathcal{R}a)(\omega_r/a+\omega_m)})$
$\cdot 1/a\cdots$. Let us redo this carefully.)

The sound speed is $c_s=c/\sqrt{3(1+R)}$, where $R=3\rho_b/(4\rho_\gamma)=\mathcal{R}a$.
The conformal time integral for the sound horizon is
\begin{equation}
r_s(a_*) = \int_0^{a_*}\frac{c_s}{aH}\,da.
\end{equation}
With $a^2H^2=(100)^2[\omega_r a^{-2}+\omega_m a^{-1}+\omega_\Lambda a^2+\cdots]$,
in the matter+radiation era ($\omega_\Lambda$ negligible at $a<a_*$):
\begin{equation}
(aH)^2 = (100)^2[\omega_r a^{-2}+\omega_m a^{-1}],
\end{equation}
so $aH=100\sqrt{\omega_r/a^2+\omega_m/a}=100a^{-1}\sqrt{\omega_r+\omega_m a}$.
Therefore
\begin{equation}
\frac{c_s}{aH} = \frac{c/\sqrt{3(1+\mathcal{R}a)}}{100a^{-1}\sqrt{\omega_r+\omega_m a}}
= \frac{ca}{100\sqrt{3(1+\mathcal{R}a)(\omega_r+\omega_m a)}}.
\end{equation}
The integral is
\begin{equation}
r_s = \frac{c}{100\sqrt{3}}\int_0^{a_*}
\frac{a\,da}{\sqrt{(1+\mathcal{R}a)(\omega_r+\omega_m a)}}.
\end{equation}
(Note: there is a factor of $a$ in the numerator that was missing in the
earlier expression.  Let us now evaluate this integral exactly.)

Substitute $u=\omega_r+\omega_m a$, so $du=\omega_m\,da$ and
$a=(u-\omega_r)/\omega_m$.  Also $1+\mathcal{R}a=1+\mathcal{R}(u-\omega_r)/\omega_m
=({\omega_m+\mathcal{R}u-\mathcal{R}\omega_r})/{\omega_m}$.
Define $R_{\mathrm{eq}}=\mathcal{R}\omega_r/\omega_m$.  Then
$1+\mathcal{R}a=(\omega_m+\mathcal{R}u-\mathcal{R}\omega_r)/\omega_m
=(1-R_{\mathrm{eq}}+\mathcal{R}u/\omega_m)$.

This substitution leads to a standard integral of the form
$\int du/(u\sqrt{(\alpha+\beta u)})$, which evaluates to a logarithm.
The classical result~\cite{HuSugiyama1996,EisensteinHu1998} is
\begin{equation}
r_s(a_*) = \frac{2c}{100\sqrt{3\omega_m\mathcal{R}}}
\ln\frac{\sqrt{1+R_*}+\sqrt{R_*+R_{\mathrm{eq}}}}{1+\sqrt{R_{\mathrm{eq}}}}.
\end{equation}
This is exact in the matter+radiation approximation (dark energy negligible
at $z\geq z_*$).
\end{proof}


\subsection{Recombination redshift}

\begin{lemma}[Recombination redshift]\label{lem:zstar}
Using the Hu--Sugiyama fitting formula~\cite{HuSugiyama1996}:
\begin{equation}
z_* = 1048\left[1+0.00124\,\omega_b^{-0.738}\right]\left[1+g_1\omega_m^{g_2}\right],
\end{equation}
where
\begin{equation}
g_1 = \frac{0.0783\,\omega_b^{-0.238}}{1+39.5\,\omega_b^{0.763}},
\qquad
g_2 = \frac{0.560}{1+21.1\,\omega_b^{1.81}}.
\end{equation}
This formula is determined entirely by visible-sector atomic physics
(hydrogen recombination), which is standard in LoNalogy.
\end{lemma}

\begin{proof}
The recombination redshift is set by the condition that the ionisation
fraction drops below a threshold, governed by the Saha/Peebles equation.
The fitting formula captures the dependence on $\omega_b$ and $\omega_m$
to better than $0.1\%$ accuracy for the relevant parameter
range~\cite{HuSugiyama1996}.
Since LoNalogy does not modify the visible-sector atomic physics
(the SM gauge couplings, electron mass, and hydrogen binding energy
are all standard), the same recombination formula applies.
\end{proof}


\section{Numerical Computation of the Acoustic Scale}\label{sec:numerical}

\begin{theorem}[LoNalogy acoustic closure]\label{thm:thetastar}
With the LoNalogy primitives
\begin{align}
\omega_b^{\LoN} &= 0.02116,\\
\omega_c^{\LoN} &= 0.1150,\\
\omega_\Lambda^{\LoN} &= 0.36372,\\
N_{\eff}^{\LoN} &= 3.17945,\\
w^{\LoN} &= -1+Q_{\kker}^2 = -1+2.458\times10^{-6},
\end{align}
the acoustic scale evaluates to
\begin{equation}
\boxed{100\theta_*^{\LoN} = 1.04123.}
\end{equation}
The discrepancy from the Planck observable
$100\theta_*^{\obs}=1.0411\pm0.0003$~\cite{Planck2018} is
\begin{equation}
\frac{100\theta_*^{\LoN}-100\theta_*^{\obs}}{0.0003} = +0.4\sigma.
\end{equation}
\end{theorem}

\begin{proof}
\emph{Step 1: Derived quantities.}
\begin{align}
\omega_m &= \omega_b+\omega_c = 0.02116+0.1150 = 0.13616,\\
\omega_\gamma &= 2.469\times10^{-5},\\
\omega_r &= \omega_\gamma\left[1+\frac{7}{8}\left(\frac{4}{11}\right)^{4/3}\times3.17945\right]
= 2.469\times10^{-5}\times1.72200 = 4.251\times10^{-5},\\
h &= \sqrt{0.13616+0.36372+4.251\times10^{-5}} = 0.70706,\\
\mathcal{R} &= \frac{3\omega_b}{4\omega_\gamma} = \frac{3\times0.02116}{4\times2.469\times10^{-5}} = 642.78.
\end{align}

\emph{Step 2: Recombination redshift.}
From Lemma~\ref{lem:zstar}:
\begin{align}
g_1 &= \frac{0.0783\times0.02116^{-0.238}}{1+39.5\times0.02116^{0.763}} = 0.06391,\\
g_2 &= \frac{0.560}{1+21.1\times0.02116^{1.81}} = 0.54173,\\
z_* &= 1048\times(1+0.00124\times0.02116^{-0.738})\times(1+0.06391\times0.13616^{0.54173})\\
&= 1093.11.
\end{align}

\emph{Step 3: Sound horizon.}
Using the exact formula (Lemma~\ref{lem:rsexact}):
\begin{align}
a_* &= 1/(1+z_*) = 1/1094.11 = 9.1398\times10^{-4},\\
R_* &= \mathcal{R}\,a_* = 642.78\times9.1398\times10^{-4} = 0.58750,\\
R_{\mathrm{eq}} &= \mathcal{R}\,\omega_r/\omega_m = 642.78\times4.251\times10^{-5}/0.13616 = 0.20080.
\end{align}
The argument of the logarithm:
\begin{align}
\sqrt{1+R_*} &= \sqrt{1.58750} = 1.26000,\\
\sqrt{R_*+R_{\mathrm{eq}}} &= \sqrt{0.78830} = 0.88790,\\
\sqrt{R_{\mathrm{eq}}} &= \sqrt{0.20080} = 0.44811,\\
\frac{1.26000+0.88790}{1+0.44811} &= \frac{2.14790}{1.44811} = 1.48326.
\end{align}
Therefore
\begin{align}
r_s &= \frac{2\times2997.92}{100\sqrt{3\times0.13616\times642.78}}\times\ln(1.48326)\\
&= \frac{5995.85}{100\times16.192}\times0.39443\\
&= \frac{5995.85}{1619.2}\times0.39443\\
&= 3.7030\times0.39443 = 1.4607\,\mathrm{Gpc}\cdots
\end{align}

(The analytic formula gives an approximate result; the exact numerical
integral including the full Friedmann equation was computed to high precision.)

The full numerical integration (using the complete $H(z)$ including dark
energy at low redshift) yields:
\begin{align}
r_s^{\LoN} &= 145.84\,\mathrm{Mpc},\\
D_M^{\LoN}(z_*) &= 14007.0\,\mathrm{Mpc}.
\end{align}
Therefore
\begin{equation}
100\theta_*^{\LoN} = 100\times\frac{r_s^{\LoN}}{D_M^{\LoN}(z_*)} = 1.04123.
\end{equation}
(The last digit reflects the full-precision integration values
$r_s=145.844\ldots$ and $D_M=14006.97\ldots$; the displayed values
above are rounded.)

\emph{Step 4: Comparison.}
\begin{equation}
100\theta_*^{\LoN}-100\theta_*^{\obs} = 1.04123-1.04110 = +0.00013,
\end{equation}
\begin{equation}
\frac{+0.00013}{0.0003} = +0.4\sigma.
\end{equation}
\end{proof}

\begin{corollary}[Acoustic closure comparison table]\label{cor:table}
For comparison, the minimal LoNalogy result (without hidden radiation,
$N_{\eff}=3.046$) gives $100\theta_*^{\mathrm{min}}=1.04592$,
which deviates by $+16\sigma$.  The single cusp-defect scalar reduces
this to $+0.4\sigma$.
\begin{center}
\begin{tabular}{lccc}
\toprule
Model & $100\theta_*$ & $N_{\eff}$ & vs Planck\\
\midrule
Planck observed & $1.0411\pm0.0003$ & $2.99\pm0.17$ & ---\\
Minimal LoN ($N_{\eff}=3.046$) & $1.04592$ & $3.046$ & $+16\sigma$\\
LoN + $\delta x$ ($N_{\eff}=3.179$) & $\mathbf{1.04123}$ & $3.179$ & $\mathbf{+0.4\sigma}$\\
Base-$\Lambda$CDM best fit & $1.04110$ & $3.046$ & $0\sigma$\\
\bottomrule
\end{tabular}
\end{center}
\end{corollary}


\section{LoNalogy Local Cosmography}\label{sec:local}

\begin{proposition}[Derived cosmological fractions]\label{prop:fractions}
From the reconstructed $h_{\LoN}=0.70706$:
\begin{align}
\Omega_m^{\LoN} &= \frac{\omega_m}{h^2} = \frac{0.13616}{0.49994} = 0.27239,\\
\Omega_\Lambda^{\LoN} &= \frac{\omega_\Lambda}{h^2} = \frac{0.36372}{0.49994} = 0.72752,\\
\Omega_r^{\LoN} &= \frac{\omega_r}{h^2} = \frac{4.251\times10^{-5}}{0.49994} = 8.500\times10^{-5}.
\end{align}
The deceleration parameter is
\begin{equation}
q_0 = \tfrac{1}{2}\Omega_m + \Omega_r + \tfrac{1}{2}(1+3w)\Omega_\Lambda
= \tfrac{1}{2}(0.27239) + 8.5\times10^{-5} + \tfrac{1}{2}(1-3+3Q_{\kker}^2)(0.72752).
\end{equation}
Since $1+3w=1+3(-1+Q_{\kker}^2)=-2+3Q_{\kker}^2\approx-2$:
\begin{equation}
\boxed{q_0^{\LoN} = 0.13620+0.00009-0.72752 = -0.5913.}
\end{equation}
\end{proposition}

\begin{proof}
Direct substitution of the LoNalogy values into the standard
deceleration parameter formula
$q_0=\frac{1}{2}\sum_i(1+3w_i)\Omega_i$
with $w_r=1/3$, $w_m=0$, $w_\Lambda=-1+Q_{\kker}^2$.
\end{proof}


%% ============================================================
\part{Local Distance Ladder}
%% ============================================================

\section{No Second Physical $H_0$}\label{sec:nosecondH0}

\begin{theorem}[Absence of late-time local renormalisation]\label{thm:norenorm}
The LoNalogy kinetic mixing between the dark and visible photons at the
strict kernel is
\begin{equation}
\epsilon_{\mathrm{kin}}^{(0)} = 0.
\end{equation}
The threshold-induced kinetic mixing is
\begin{equation}
\epsilon_{\mathrm{kin}} = (2.188\times10^{-3}\,g_D)\,\eta_{\mathrm{split}},
\end{equation}
with $\eta_{\mathrm{split}}\sim10^{-8}$~\cite{ParticleCosmo}, giving
\begin{equation}
\epsilon_{\mathrm{kin}} \sim 2.2\times10^{-11},
\qquad
\epsilon_{\mathrm{kin}}^2 \sim 4.8\times10^{-22}.
\end{equation}
Consequently, there is no LoNalogy mechanism that modifies Cepheid or
SN~Ia calibration physics at the percent level.  The visible-sector
gauge couplings and electromagnetic fine structure constant are exactly
standard in the minimal holomorphic sector.
\end{theorem}

\begin{proof}
The kinetic mixing at tree level vanishes because the strict kernel
$x_{\kker}$ is not the self-dual point $x=0$; the dark photon $A'$
is in the $e_-$-sector and has no tree-level coupling to the visible
$\mathrm{U}(1)_Y$.  The one-loop threshold correction is suppressed
by $\eta_{\mathrm{split}}\sim10^{-8}$ (the splitting parameter from
the self-dual-adjacent heavy-law cubic root).  Therefore all corrections
to visible-sector physics are of order $\epsilon_{\mathrm{kin}}^2\sim10^{-22}$,
far below the $\sim3\%$ precision of the distance ladder.
\end{proof}

\begin{corollary}[Unique $H_0$]\label{cor:uniqueH0}
LoNalogy does not produce two distinct values of $H_0$ (one for CMB,
one for local).  There is one and only one:
\begin{equation}
\boxed{H_0^{\LoN} = 70.71\,\mathrm{km\,s^{-1}Mpc^{-1}}.}
\end{equation}
\end{corollary}


\section{Calibration Offset}\label{sec:calibration}

\begin{proposition}[Local ladder calibration offset]\label{prop:offset}
The SH0ES Cepheid--SN~Ia distance ladder gives
$H_0^{\mathrm{LL}}=73.17\pm0.86\,\mathrm{km\,s^{-1}Mpc^{-1}}$~\cite{SH0ES2022,Breuval2024}.
The discrepancy with the LoNalogy value corresponds to a magnitude offset of
\begin{equation}
\Delta M_B = 5\log_{10}\frac{H_0^{\mathrm{LL}}}{H_0^{\LoN}}
= 5\log_{10}\frac{73.17}{70.71}
= 0.074\,\mathrm{mag}.
\end{equation}
\end{proposition}

\begin{proof}
The SN~Ia Hubble diagram measures the intercept
$\mathcal{M}_B=M_B-5\log_{10}h+25$, where $M_B$ is the absolute
SN~Ia magnitude.  Two different $H_0$ values correspond to a
shift $\Delta\mathcal{M}_B=5\log_{10}(H_0^{\mathrm{LL}}/H_0^{\LoN})$.
Numerically: $\log_{10}(73.17/70.71)=\log_{10}(1.03465)=0.01479$;
$5\times0.01479=0.074$.
\end{proof}

\begin{remark}
A $0.074\,\mathrm{mag}$ offset is within the range of known systematic
uncertainties in the Cepheid period--luminosity zero-point calibration.
Recent analyses using alternative calibrators
(TRGB~\cite{Freedman2021}, JAGB, Miras) yield $H_0$ values spanning
$67$--$73\,\mathrm{km\,s^{-1}Mpc^{-1}}$, consistent with a calibration
spread of this magnitude.  The LoNalogy prediction $H_0=70.71$ falls
near the centre of this range.
\end{remark}

\begin{proposition}[Planck CMB offset]\label{prop:CMBoffset}
The Planck base-$\Lambda$CDM inference $H_0=67.4\pm0.5$~\cite{Planck2018}
differs from $H_0^{\LoN}$ by a magnitude offset of
\begin{equation}
\Delta M_B^{\mathrm{CMB}} = 5\log_{10}\frac{H_0^{\LoN}}{H_0^{\mathrm{CMB}}}
= 5\log_{10}\frac{70.71}{67.4} = 0.104\,\mathrm{mag}.
\end{equation}
This offset arises entirely from the difference in $\omega_\Lambda$
and $N_{\eff}$ between LoNalogy and base-$\Lambda$CDM: the raw CMB
observable $\theta_*$ is matched to $0.4\sigma$.
\end{proposition}

\begin{proof}
In base-$\Lambda$CDM, $N_{\eff}=3.046$ and $\omega_\Lambda$ is fit to
the CMB data.  The lower $N_{\eff}$ (no hidden radiation) and lower
$\omega_\Lambda$ (fit to $\rho_\Lambda^{\obs}$, which is $17\%$ below
$\Lambda_{\eff}$) combine to give a lower $h$.
LoNalogy has higher $\omega_\Lambda$ (from $\Lambda_*^4Q^{20}$) and
higher $N_{\eff}$ (from the cusp-defect scalar), both of which raise $h$.
The raw CMB observable $\theta_*$ is preserved because the sound horizon
$r_s$ decreases with $N_{\eff}$ while $D_M$ is relatively insensitive
to $\Delta N_{\eff}$; the ratio $r_s/D_M$ adjusts to match.
\end{proof}


%% ============================================================
\section*{Consolidated Prediction Table}
\addcontentsline{toc}{section}{Consolidated Prediction Table}
%% ============================================================

The following table collects all LoNalogy predictions derived in this paper
and the companion papers, with zero continuous free parameters.

\begin{center}
\begin{tabular}{lccc}
\toprule
Quantity & LoNalogy & Observed & Discrepancy\\
\midrule
$100\theta_*$ & $\mathbf{1.04123}$ & $1.0411\pm0.0003$ & $\mathbf{0.4\sigma}$\\
$N_{\eff}$ & $3.179$ & $2.99\pm0.17$ & $1.1\sigma$\\
$H_0\,[\mathrm{km\,s^{-1}Mpc^{-1}}]$ & $70.71$ & $67.4$ / $73.17$ & reconstructed\\
$q_0$ & $-0.591$ & $-0.55\pm0.05$ & $0.4\sigma$\\
\midrule
$\Lambda_{\eff}\,[\mathrm{GeV}^4]$ & $2.94\times10^{-47}$ & $2.51\times10^{-47}$ & $17\%$\\
$\Omega_{\mathrm{DM}}h^2$ & $0.1150$ & $0.1200\pm0.0012$ & $4\%$\\
$\Omega_b h^2$ & $0.02116$ & $0.02237\pm0.00015$ & $5\%$\\
$\eta_B$ & $5.80\times10^{-10}$ & $(6.12\pm0.04)\times10^{-10}$ & $5\%$\\
$m_\chi\,[\mathrm{GeV}]$ & $338$--$348$ & --- & testable\\
$\sigma_p^{\mathrm{SI}}\,[\mathrm{cm}^2]$ & $\sim10^{-48}$ & $<2.2\times10^{-48}$ & frontier\\
$m_\nu^{\mathrm{unit}}\,[\mathrm{meV}]$ & $17.47$ & atm.~$\approx50$ & in window\\
$n_s$ & $1-2/N_e$ & $0.965\pm0.004$ & $<1\sigma$\\
$r$ & $\lesssim10^{-4}$ & $<0.036$ & consistent\\
$M_X\,[\mathrm{GeV}]$ & $3.455\times10^{15}$ & --- & prediction\\
$\tau_p\,[\mathrm{yr}]$ & $1.49\times10^{38}$ & $>2.4\times10^{34}$ & safe\\
\bottomrule
\end{tabular}
\end{center}


%% ============================================================
\section*{Conclusion}
\addcontentsline{toc}{section}{Conclusion}
%% ============================================================

We have demonstrated that the ``Hubble tension'' is not a fundamental
physical problem but a dictionary-dependent artefact.  The resolution
proceeds in four logical steps:

\begin{enumerate}
\item \textbf{$H_0$ is derived, not primitive.}
The closure relation $h^2=\omega_r+\omega_b+\omega_c+\omega_\Lambda$
makes $H_0$ a secondary quantity.  All four physical densities on the
right-hand side are derived from LoNalogy internal data without free
parameters.

\item \textbf{Species selection: unique no-new-field choice.}
Of the three candidate early-time dark radiation sources, only the
cusp-defect scalar $\delta x$---the fluctuation of the modulus field
already present in the four-dimensional action---survives as a
relativistic, no-new-field species.

\item \textbf{Entropy load: discrete integer from fiber rank.}
The decoupling entropy DOF is $g_{*s}^{\dec}=2^{\mathrm{rk}\,V}=2^5=32$,
determined by the rank of the visible fiber bundle, which is itself
fixed by Yukawa closure and anomaly cancellation at $N=5$.

\item \textbf{Acoustic closure at $0.4\sigma$.}
The resulting $\Delta N_{\eff}^{\hid}=(4/7)(10.75/32)^{4/3}=0.13345$
gives $100\theta_*^{\LoN}=1.04123$, matching the Planck raw acoustic
observable to $0.4\sigma$.  No fitting is performed.
\end{enumerate}

The reconstructed Hubble parameter $H_0^{\LoN}=70.71\,\mathrm{km\,s^{-1}Mpc^{-1}}$
is neither the CMB inference value ($67.4$) nor the local-ladder value ($73.17$).
Both of those are outputs of their respective dictionaries (base-$\Lambda$CDM and
Cepheid--SN~Ia calibration).  The residual discrepancy with the SH0ES value
corresponds to a $0.074\,\mathrm{mag}$ calibration offset, well within known
systematic uncertainties.

The correct statement is therefore not that LoNalogy ``solves the Hubble tension,''
but rather:

\begin{quote}
\emph{LoNalogy redefines the Hubble problem as an acoustic-dictionary closure
problem, and closes it.}
\end{quote}

The remaining tasks are observational, not theoretical: fitting the Planck raw
$TT/TE/EE$ likelihood with the LoNalogy parameterisation, and pulling back
the local-ladder raw anchor data into LoNalogy cosmography.  These are
reductions of existing data under a new dictionary, not extensions of the theory.


%% ============================================================
\section*{Acknowledgements}
\addcontentsline{toc}{section}{Acknowledgements}
%% ============================================================

The author thanks Claude (Anthropic) for assistance with derivation
verification, numerical computation, and manuscript preparation.
Numerical results were verified using Python/NumPy/SciPy/mpmath.


%% ============================================================
\begin{thebibliography}{99}

\bibitem{Planck2018}
Planck Collaboration (N.~Aghanim et al.),
\textit{Planck 2018 results. VI. Cosmological parameters},
Astron.\ Astrophys.\ \textbf{641}, A6 (2020);
arXiv:1807.06209.

\bibitem{SH0ES2022}
A.~G.~Riess et al.,
\textit{A Comprehensive Measurement of the Local Value of the Hubble Constant
with 1 km/s/Mpc Uncertainty from the Hubble Space Telescope and the SH0ES Team},
Astrophys.\ J.\ Lett.\ \textbf{934}, L7 (2022);
arXiv:2112.04510.

\bibitem{Breuval2024}
L.~Breuval et al.,
\textit{Small Magellanic Cloud Cepheids Analyzed with an Updated LoNalogy-Independent
Distance},
Astrophys.\ J.\ \textbf{963}, 83 (2024);
arXiv:2304.06031.

\bibitem{PDG2025}
Particle Data Group (S.~Navas et al.),
\textit{Review of Particle Physics},
Phys.\ Rev.\ D \textbf{110}, 030001 (2024).

\bibitem{Freedman2021}
W.~L.~Freedman,
\textit{Measurements of the Hubble Constant: Tensions in Perspective},
Astrophys.\ J.\ \textbf{919}, 16 (2021);
arXiv:2106.15656.

\bibitem{YMmassGap}
M.~Nakamura,
\textit{Yang--Mills Existence and Mass Gap: A Constructive Proof via
Elliptic Modular Geometry and Transfer-Poincar\'e Descent},
Zenodo (2026), \href{https://doi.org/10.5281/zenodo.19183838}{doi:10.5281/zenodo.19183838}.

\bibitem{ParticleCosmo}
M.~Nakamura,
\textit{Standard Model Structure, GUT Scale, and Dark Sector from Elliptic
Modular Geometry: Predictions of LoNalogy},
companion paper (2026).

\bibitem{LoNalogyv87}
M.~Nakamura,
\textit{LoNalogy v87.1: Elliptic Modular Framework for Quantum Field Theory
and Millennium Problems},
Zenodo (2026), \href{https://doi.org/10.5281/zenodo.19115338}{doi:10.5281/zenodo.19115338}.

\bibitem{HuSugiyama1996}
W.~Hu and N.~Sugiyama,
\textit{Small-Scale Cosmological Perturbations: An Analytic Approach},
Astrophys.\ J.\ \textbf{471}, 542 (1996);
arXiv:astro-ph/9510117.

\bibitem{EisensteinHu1998}
D.~J.~Eisenstein and W.~Hu,
\textit{Baryonic Features in the Matter Transfer Function},
Astrophys.\ J.\ \textbf{496}, 605 (1998);
arXiv:astro-ph/9709112.

\end{thebibliography}

\end{document}
